% F081: precise mathematical teaching diagram.
\begin{tikzpicture}[x=1cm,y=1cm,>=stealth,font=\small]
\draw[rounded corners=3pt,gray!40] (0,0) rectangle (14.2,9.4);
\begin{scope}[yshift=1cm]
\node at (7.1,7.800000000000001) {易错二：导数零点一样多，原函数零点未必一样多};
\node[] at (3.7,7) {$x^3-3x+1$：3个零点};
\begin{scope}[shift={(3.7,3.2)},x=0.92cm,y=0.42cm]
\draw[->,gray!65] (-2.9,0)--(2.6,0) node[right] {$x$};
\draw[->,gray!65] (0,-3.8)--(0,7.2) node[left] {$y$};
\begin{scope}\clip (-2.7,-3.5) rectangle (2.4,7.7);
\draw[blue!75!black,line width=1.15pt,,domain=-2.7:2.4,samples=160] plot (\x,{(\x)*(\x)*(\x)-3*\x+1});
\end{scope}
\fill[blue!75!black] (-1.8793852415718,0) circle (1.8pt);
\fill[blue!75!black] (0.3472963553339,0) circle (1.8pt);
\fill[blue!75!black] (1.532088886238,0) circle (1.8pt);
\draw[orange!85!black,line width=1.2pt] (-1.3,3)--(-0.7,3);
\fill[orange!85!black] (-1,3) circle (1.8pt);
\draw[orange!85!black,line width=1.2pt] (0.7,-1)--(1.3,-1);
\fill[orange!85!black] (1,-1) circle (1.8pt);
\end{scope}
\node[] at (10.4,7) {$x^3-3x+5$：1个零点};
\begin{scope}[shift={(10.4,3.2)},x=0.92cm,y=0.42cm]
\draw[->,gray!65] (-2.9,0)--(2.6,0) node[right] {$x$};
\draw[->,gray!65] (0,-3.8)--(0,7.2) node[left] {$y$};
\begin{scope}\clip (-2.7,-3.5) rectangle (2.4,7.7);
\draw[blue!75!black,line width=1.15pt,,domain=-2.7:2.4,samples=160] plot (\x,{(\x)*(\x)*(\x)-3*\x+5});
\end{scope}
\fill[blue!75!black] (-2.2790187861666,0) circle (1.8pt);
\draw[orange!85!black,line width=1.2pt] (-1.3,7)--(-0.7,7);
\fill[orange!85!black] (-1,7) circle (1.8pt);
\draw[orange!85!black,line width=1.2pt] (0.7,3)--(1.3,3);
\fill[orange!85!black] (1,3) circle (1.8pt);
\end{scope}
\end{scope}
\node[text width=13cm,align=center] at (7.1,.9) {两图同尺度，导数都只有两个零点；只能说最多三个，不能说恰好三个。};
\end{tikzpicture}
